\chapter{Prostate-Specific Antigen (PSA)}

\section*{What Is This Test?}
Prostate-specific antigen (PSA) is a glycoprotein enzyme (serine protease, also known as kallikrein-3 or hK3) produced almost exclusively by the epithelial cells of the prostate gland. PSA functions in semen to liquefy the seminal coagulum after ejaculation by cleaving semenogelin I and II and fibronectin, allowing spermatozoa to become motile. A small amount of PSA normally leaks into the bloodstream, where it circulates in two forms: complexed PSA (bound to protease inhibitors, primarily alpha-1-antichymotrypsin and alpha-2-macroglobulin) and free PSA (unbound, enzymatically inactive). Total PSA (tPSA) is the sum of free PSA (fPSA) and complexed PSA. In prostate cancer, disruption of the normal prostatic architecture allows more PSA to enter the circulation, leading to elevated serum levels. PSA testing is used for prostate cancer screening (controversially), diagnosis, risk stratification, monitoring of treatment response, and detection of biochemical recurrence. The test is performed by immunoassay (electrochemiluminescence, chemiluminescence, or ELISA) using monoclonal antibodies that recognize distinct epitopes on the PSA molecule.

\section*{Alternative Names}
PSA; Total PSA (tPSA); Free PSA (fPSA); Percent Free PSA (\%fPSA); Prostate-Specific Antigen; Kallikrein-3; hK3; PSA Density; PSA Velocity; PSA Doubling Time

\section*{Principle}
PSA is measured using a two-site sandwich electrochemiluminescent immunoassay (ECLIA) or chemiluminescent immunoassay (CLIA) on automated immunoassay analyzers (Roche Cobas, Abbott Architect, Siemens Atellica). In the ECLIA method, patient serum is incubated with a biotinylated monoclonal anti-PSA capture antibody and a ruthenium-complex-labeled monoclonal anti-PSA detection antibody, forming a sandwich immunocomplex. Streptavidin-coated paramagnetic microparticles capture the biotinylated antibody. The reaction mixture is aspirated into the measuring cell, where microparticles are magnetically captured onto an electrode. Unbound reagents are washed away, and tripropylamine (TPA) is added. An electrical voltage applied to the electrode induces a chemiluminescent reaction between the ruthenium complex and TPA, emitting light at 620 nm that is measured by a photomultiplier. The light signal is proportional to PSA concentration. Total PSA and free PSA are measured in separate assays using antibodies that recognize different epitopes. Equimolar detection of free and complexed PSA is essential for accurate total PSA measurement. PSA calibrators are traceable to WHO International Standard (96/670 for tPSA, 96/668 for fPSA). Percent free PSA is calculated as (fPSA/tPSA) $\times$ 100.

\section*{History}
PSA was discovered in 1970 by Richard Ablin and colleagues in seminal plasma, originally as a forensic marker for semen identification. The purification and characterization of PSA as a prostate-specific protein was performed by Ming Wang and T.M. Chu in 1979 at Roswell Park Memorial Institute. The first serum PSA immunoassay was developed by the same group in 1980. The FDA approved the PSA test for monitoring prostate cancer recurrence in 1986 and for screening (in conjunction with digital rectal examination, DRE) in 1994. The discovery that free-to-total PSA ratio improves specificity for prostate cancer detection in the diagnostic gray zone (PSA 4--10 ng/mL) was made by Catalona et al. and Christensson et al. in the 1990s. Large randomized screening trials---the Prostate, Lung, Colorectal and Ovarian (PLCO) Cancer Screening Trial in the United States and the European Randomized Study of Screening for Prostate Cancer (ERSPC)---published in 2009--2012, demonstrated that PSA screening reduces prostate cancer mortality by approximately 20--30\% but with a high rate of overdiagnosis and overtreatment. The 2012 U.S. Preventive Services Task Force (USPSTF) recommended against routine PSA screening (Grade D), later revised in 2018 to recommend shared decision-making for men aged 55--69 (Grade C) \cite{uspstf2018prostate}. The introduction of multiparametric MRI (mpMRI) and MRI-targeted biopsy has improved risk stratification before biopsy, reducing unnecessary biopsies and detection of clinically insignificant prostate cancer \cite{kasivisvanathan2018mri}.

\section*{Why This Test Is Ordered}
\begin{itemize}
  \item Screening for prostate cancer in men aged 55--69 as part of shared decision-making, per USPSTF (2018) and American Urological Association (AUA) guidelines. Screening is considered in men aged 40--54 with high risk: African American men, men with first-degree relatives with prostate cancer (especially at a young age), and men with BRCA1/BRCA2 or Lynch syndrome mutations
  \item Diagnosis and risk stratification of prostate cancer in men with elevated PSA or abnormal digital rectal examination (DRE)
  \item Monitoring for treatment response and detection of biochemical recurrence after radical prostatectomy (PSA should be undetectable, as all prostate tissue has been removed), external beam radiation therapy, or brachytherapy
  \item Active surveillance of low-risk prostate cancer --- serial PSA measurement to detect disease progression
  \item Percent free PSA (\%fPSA) measurement when total PSA is in the diagnostic gray zone (4--10 ng/mL) to improve specificity before proceeding to prostate biopsy. A high \%fPSA (>25\%) suggests benign prostatic hyperplasia; a low \%fPSA (<10\%) is more suspicious for prostate cancer \cite{catalona1998percentage}
  \item PSA doubling time and PSA velocity (rate of PSA change over time) for prognostic assessment and treatment decisions
  \item PSA monitoring in men on 5-alpha reductase inhibitors (finasteride, dutasteride) for BPH to correct for drug-related PSA reduction (PSA is multiplied by 2)
\end{itemize}

\section*{Primary vs Secondary Test}
This is a primary screening test for prostate cancer, although its role in screening is controversial due to high rates of false-positive results and overdiagnosis. PSA is also a primary test for monitoring prostate cancer after treatment.

\section*{How This Test Is Performed}
For most patients, a health care professional will take a blood sample from a vein in your arm using a small needle. After the needle is inserted, a small amount of blood will be collected into a test tube or vial. You may feel a little sting when the needle goes in or out. This usually takes less than five minutes.

\section*{What Preparation Is Needed}
Critical pre-analytical instructions: PSA should NOT be drawn after digital rectal examination (DRE), prostate biopsy, cystoscopy, transrectal ultrasound, ejaculation (avoid ejaculation for 48 hours before testing), or vigorous bicycle riding, as these activities mechanically disrupt the prostatic epithelium and cause transient PSA elevation. PSA should be drawn BEFORE a DRE if both are planned during the same visit. PSA sampling should be delayed for at least 4--6 weeks after acute prostatitis, urinary tract infection, or acute urinary retention. No fasting is required. Patients should inform their healthcare provider if they are taking 5-alpha reductase inhibitors (finasteride, dutasteride), which lower PSA by approximately 50\% after 6--12 months of therapy. Chemotherapy, androgen deprivation therapy, and radiation therapy will lower PSA levels.

\section*{Contraindications and Precautions}
No absolute contraindications to venipuncture. The primary limitation of PSA testing is its lack of specificity --- PSA is organ-specific (prostate) but not cancer-specific. Elevated PSA occurs in numerous non-malignant conditions: benign prostatic hyperplasia (BPH, the most common cause of mild to moderate PSA elevation in older men), acute and chronic prostatitis, acute urinary retention, recent instrumentation (DRE, cystoscopy, urethral catheterization, prostate biopsy), and ejaculation. Approximately 75\% of men with PSA 4--10 ng/mL who undergo transrectal ultrasound-guided biopsy have negative biopsies (false-positive PSA). Conversely, prostate cancer can exist with a PSA <4.0 ng/mL --- the Prostate Cancer Prevention Trial (PCPT) found that 15.2\% of men with PSA <4.0 ng/mL and a normal DRE had prostate cancer on biopsy, including 2.3\% with high-grade (Gleason $\geq$7) cancer \cite{thompson2004prevalence}. PSA levels increase with age due to age-related prostatic enlargement. Certain common laboratory interferences: heterophile antibodies and human anti-mouse antibodies (HAMA) can cause falsely elevated PSA results. Biotin (vitamin B7) supplementation at high doses (>5 mg/day) can interfere with biotin-streptavidin immunoassay systems. Specimen should be collected in a serum separator tube (SST, gold-top) or plasma (lithium heparin, green-top tube). Serial PSA measurements should ideally be performed using the same PSA assay methodology (same platform) for consistency.

\section*{Risks}
Minimal risk from venipuncture. The major risk associated with PSA testing is not from the blood draw itself but from the cascade of downstream consequences following a false-positive result: unnecessary prostate biopsies (which carry risks of infection, bleeding, pain, urinary retention, and sepsis), overdiagnosis of clinically insignificant prostate cancer (Gleason 6, low volume, unlikely to cause morbidity or mortality), and overtreatment with radical prostatectomy or radiation with attendant risks of erectile dysfunction (30--70\%), urinary incontinence (5--20\%), and bowel dysfunction.

\section*{What Happens After the Test}
Results are typically available within 1 to 2 hours. An elevated PSA should be confirmed by repeat testing after 4--6 weeks, particularly if a transient cause of elevation (ejaculation, recent instrumentation, prostatitis) cannot be excluded. The decision to proceed to prostate biopsy after a confirmed elevated PSA depends on: total PSA level, percent free PSA, PSA density (PSA/prostate volume measured by transrectal ultrasound or MRI), DRE findings, patient age, comorbidities, life expectancy, family history, and patient preference \cite{heidenreich2014eau}. Patients with elevated PSA increasingly undergo multiparametric prostate MRI before biopsy: a negative MRI (PI-RADS score 1--2) has a high negative predictive value for clinically significant prostate cancer and may obviate biopsy. An indeterminate or positive MRI (PI-RADS 3--5) guides targeted biopsy of suspicious lesions.

\section*{Understanding Results}

\subsection*{Below Normal}
\begin{itemize}
  \item PSA <4.0 ng/mL: Considered within the age-adjusted reference range for most men. Does not exclude prostate cancer; approximately 15\% of men with PSA <4.0 ng/mL and a normal DRE have prostate cancer on biopsy (PCPT trial)
  \item Undetectable PSA (<0.1 ng/mL) after radical prostatectomy: Indicates complete surgical removal of prostate tissue (clinical and biochemical remission). A subsequent detectable and rising PSA (>0.2 ng/mL confirmed) defines biochemical recurrence and warrants evaluation for local recurrence or distant metastases
  \item PSA nadir <0.5 ng/mL after radiation therapy (EBRT or brachytherapy): Indicates favorable treatment response. PSA does not become undetectable after radiation because benign prostate tissue remains and continues to produce a small amount of PSA. Biochemical failure after radiation is defined by the Phoenix criteria (nadir + 2 ng/mL rise)
  \item PSA reduction of approximately 50\% in men taking 5-alpha reductase inhibitors (finasteride, dutasteride) for BPH: Expected physiologic response due to regression of benign prostatic epithelium. The PSA value should be multiplied by 2 when interpreting PSA in men taking these medications chronically (>6 months)
  \item Androgen deprivation therapy (ADT) with GnRH agonists (leuprolide, goserelin) or antagonists (degarelix) reduces PSA to castrate levels (<0.2 ng/mL). Rising PSA on ADT defines castration-resistant prostate cancer (CRPC)
\end{itemize}

\subsection*{Above Normal}
\begin{itemize}
  \item PSA 4--10 ng/mL (diagnostic gray zone): The most common range for PSA in clinical practice. Prostate cancer is found on biopsy in approximately 25\% of men in this range. Percent free PSA helps refine the risk: \%fPSA >25\% suggests BPH (high negative predictive value), \%fPSA <10\% is more suspicious for prostate cancer and prompts biopsy
  \item PSA >10 ng/mL: Prostate cancer is found in >50\% of men undergoing biopsy. This level warrants strong consideration for prostate biopsy regardless of DRE findings
  \item PSA >20 ng/mL: Associated with a high risk of extracapsular extension and seminal vesicle invasion. Staging (CT abdomen/pelvis, bone scan, or PSMA-PET/CT) should be performed
  \item PSA >100 ng/mL: Almost always indicates disseminated metastatic prostate cancer, though an isolated very high PSA can occasionally be due to severe acute prostatitis
  \item PSA velocity: PSA change over time may contribute to risk assessment, but velocity should not be used as a stand-alone indication for biopsy. Use repeat PSA, DRE, risk calculators, MRI, family history, life expectancy, and shared decision-making according to current guidance.
  \item PSA doubling time <3 months: Indicates rapidly progressive, aggressive castration-resistant prostate cancer with poor prognosis
  \item False-positive PSA elevation: Acute prostatitis (PSA can rise to >50--100 ng/mL and takes 4--6 weeks to normalize), urinary tract infection, acute urinary retention, recent ejaculation, prostate biopsy, transurethral resection of the prostate (TURP), and vigorous bicycle riding
\end{itemize}

\section*{Normal Results}
\begin{itemize}
  \item \textbf{Total PSA (tPSA) --- general population screening cutoffs:}
  \begin{itemize}
    \item <4.0 ng/mL: Normal (but does NOT exclude prostate cancer; 15\% of men with PSA <4.0 ng/mL have prostate cancer)
    \item 4.0--10.0 ng/mL: Gray zone (approximately 25\% probability of prostate cancer)
    \item >10.0 ng/mL: Elevated (approximately 50--67\% probability of prostate cancer)
  \end{itemize}
  \item \textbf{Age-specific PSA reference ranges (proposed by Oesterling et al.):}
  \begin{itemize}
    \item Age 40--49: <2.5 ng/mL
    \item Age 50--59: <3.5 ng/mL
    \item Age 60--69: <4.5 ng/mL
    \item Age 70--79: <6.5 ng/mL
  \end{itemize}
  \item \textbf{Percent free PSA (\%fPSA) --- when tPSA is 4--10 ng/mL:}
  \begin{itemize}
    \item >25\%: Low probability of prostate cancer (suggests BPH)
    \item 10--25\%: Intermediate risk
    \item <10\%: Higher probability of prostate cancer
  \end{itemize}
  \item \textbf{PSA after radical prostatectomy:} Undetectable (<0.1 ng/mL). Two consecutive PSA values $\geq$0.2 ng/mL define biochemical recurrence
  \item \textbf{PSA after radiation therapy:} Nadir + 2 ng/mL defines biochemical failure (Phoenix criteria)
  \item \textbf{Men on 5-alpha reductase inhibitors ($\geq$6 months):} PSA should be multiplied by 2 for interpretation
\end{itemize}

\section*{Follow-up Tests}
\begin{itemize}
  \item \textbf{Digital rectal examination (DRE):} Palpation of the prostate through the rectal wall to detect asymmetry, nodules, induration, or frank masses. Should be performed before PSA blood draw. The combination of DRE and PSA improves detection of clinically significant prostate cancer
  \item \textbf{Percent free PSA (\%fPSA):} Automatically reflexed or separately ordered when total PSA is 4--10 ng/mL. A low \%fPSA increases suspicion for prostate cancer and reduces unnecessary biopsies by 20--30\%
  \item \textbf{Prostate Health Index (PHI):} A multivariable test that combines tPSA, fPSA, and [-2]proPSA (p2PSA, a precursor of PSA) to improve specificity for prostate cancer detection, particularly clinically significant (Gleason $\geq$7) cancer. PHI is FDA-approved for men with PSA 4--10 ng/mL and a negative DRE
  \item \textbf{4Kscore Test:} A blood test that measures tPSA, fPSA, intact PSA (iPSA), and human kallikrein 2 (hK2), combined with clinical information (age, DRE findings, prior biopsy status) to predict the risk of high-grade (Gleason $\geq$7) prostate cancer on biopsy
  \item \textbf{Multiparametric MRI (mpMRI) of the prostate:} Provides anatomical and functional imaging of the prostate. PI-RADS (Prostate Imaging Reporting and Data System) score: 1--2 (very low/low likelihood of clinically significant cancer), 3 (intermediate), 4--5 (high/very high). A negative MRI (PI-RADS 1--2) has a high negative predictive value and can avoid unnecessary biopsy
  \item \textbf{Transrectal ultrasound (TRUS)-guided prostate biopsy:} 12-core systematic biopsy with additional targeted cores from MRI-visible lesions (MRI-ultrasound fusion biopsy). The gold standard for histologic diagnosis of prostate cancer
  \item \textbf{PSMA-PET/CT (prostate-specific membrane antigen positron emission tomography):} Detects prostate cancer metastases with high sensitivity and specificity. Superior to CT and bone scan for staging high-risk and recurrent prostate cancer
  \item \textbf{CT abdomen/pelvis and technetium-99m bone scan:} For staging of newly diagnosed intermediate- or high-risk prostate cancer. Detect lymph node and bone metastases, respectively
\end{itemize}

\section*{References}
\bibliographystyle{unsrtnat}
\bibliography{references}

\clearpage
